\section{Butterfly layer and interchanges with general parameters}\label{app:parameter-transfer}\label{app:certificates}

Starting from the elementary tape operations (Fact~\ref{fact:oai-streams}, \cite[Lemmas~3--4]{OAI26}) and the finite
linear network certificates (Fact~\ref{fact:oai-networks}, \cite[Props.~7--8]{OAI26}), this appendix establishes the
exact network-exponent and guard certificates at $h=100$ (Proposition~\ref{prop:deficits}), the parameter-exposed
address interchanges (Proposition~\ref{prop:transfer11} and Lemma~\ref{lem:transfer12}), and the packed selected-bit
addition (Lemma~\ref{lem:transfer13}). From these it derives the recursive complex phase-frame reduction
(Lemma~\ref{lem:transfer-phase}) and the simultaneous butterfly layer with general structural and guard parameters
(Proposition~\ref{prop:prop14-param}). In particular, we do not invoke~\cite[Prop.~14]{OAI26} outside its printed parameter
tuple; instead we follow the written arguments of~\cite[\S\S4--5]{OAI26} step by step with all exponents and guard
parameters exposed. The finite-network constructions themselves are those of~\cite[\S\S3.1--3.3, Lemmas~5--6]{OAI26},
used here solely through the finite algebraic data of Definition~\ref{def:fixed-network-premises}
(\cite[Props.~7--8, \S\S3.4--3.5]{OAI26}).

\paragraph{Conventions.} We use the machine model and array conventions of Section~\ref{sec:prelim-model} and
Appendix~\ref{app:oai} (\cite[\S2]{OAI26}). Arrays are stored in lexicographic order with the first field most
significant, and addresses are implicit. A trailing payload of $B$ bits may be viewed as a final field of range $[B]$,
and an address permutation $f$ moves the entry at $x$ to $f(x)$. The logical volume $V$ is the number of data bits of
the current array; temporary copies and parked streams are charged in time but never replace $V$ in a recurrence. Every
machine has a fixed finite alphabet and a fixed finite number of tapes; implied constants may depend on the fixed
networks, the fixed rational exponents and the fixed comparison constants, but not on field lengths or data. Every
supplied shape is given by a fixed number of canonical binary lengths unless a longer list is explicitly charged to the
caller, and products of adjacent spectator lengths are formed before a primitive is called. As in
Sections~\ref{sec:walsh} and~\ref{sec:direct-fft}, $q$ denotes the dyadic grid precision at which the layer operates,
and $Q_q$ truncates real and imaginary components toward zero to $2^{-q}\Z$ (so that $\|Q_qz\|_\infty\le\|z\|_\infty$
and $\|Q_qz-z\|_\infty<\sqrt2\cdot2^{-q}$). Finally, $\Gamma$ is the self-delimiting header code~\eqref{eq:oai-gamma}
(\cite[(1)]{OAI26}), and $q_{\mathrm{rad}}$ is a fixed prime radix.

\subsection{Fixed network data, certified exponents, and elementary procedures}\label{sec:transfer-premises}

\begin{definition}[Fixed network data]\label{def:fixed-network-premises}
The following finite data are fixed by Fact~\ref{fact:oai-networks}:
\begin{enumerate}[label=(\alph*),leftmargin=*]
\item \emph{Bit network.} A finite pointwise XOR circuit on $W_b$ roles, auxiliary roles included, whose completed
  scalar action is a permutation $\rho_b$ of all inputs, for arbitrary initial auxiliary values. Each terminal and gate
  carries a rational $m\times m$ matrix, common to all incidences of a gate. With
  $A_e=M_{\mathrm{head}(e)}-M_{\mathrm{tail}(e)}$ and $s_b=\sum_e\operatorname{rank}_{\Q}A_e$,
  \[
    M_{\mathrm{out}(\rho_b(w))}-M_{\mathrm{in}(w)}=I_m\quad\text{for every input role }w,\qquad s_b<W_bm .
  \]
\item \emph{Complex network.} A finite scalar linear circuit on $W_c$ roles with coefficients in
  $\{0,\pm1,\pm\frac12\}$ and binary subspace labels as in~\cite[Prop.~8, \S3.5]{OAI26}. Adjacent labels are nested,
  the nonzero edge residuals have fixed orthonormal bases, and their total dimension is $s_c<W_cm$. The scalar circuit
  exchanges its two data banks with the prescribed signs and restores its auxiliary values for arbitrary initial
  values; its source, sink and bank corrections are those of~\cite[proof of Prop.~8, \S3.5, pp.~19--20]{OAI26}. We use
  only that this scalar circuit is fixed and finite; no bound on its number of gates is assumed
  (Remark~\ref{rem:transfer-guard-explicit}).
\end{enumerate}
For the triple-network constructions at $h=100$ (\cite[\S3.1, (3)--(4); Props.~7--8, \S\S3.4--3.5; \S5.4, p.~35]{OAI26}),
\[
  m=10^6,\quad W_b=177\,176\,569\,091\,445\,000\,000,\quad s_b=177\,176\,569\,088\,785\,861\,287\,000\,000,
\]
\[
  W_c=1\,873\,807\,244\,643\,542\,670\,000,\quad s_c=1\,873\,807\,244\,636\,671\,267\,308\,000\,000,
\]
with exact relative deficits (\cite[\S3.6]{OAI26})
\[
  \eta_b=\frac{W_bm-s_b}{W_bm}=\frac{339}{22\,587\,335\,000\,000},\qquad
  \eta_c=\frac{W_cm-s_c}{W_cm}=\frac{73}{19\,906\,842\,167\,500}.
\]
\end{definition}

Following~\cite[\S5.4]{OAI26}, we define from the complex-network counts $(m,W_c,s_c)$ the guard constants
\begin{equation}\label{eq:walsh-guard-constants}
  E_{\rm net}=64(W_c+m+1)^3,\quad B_{\rm net}=s_c+E_{\rm net},\quad
  \nu_{\mathrm{net}}=\min\{j\ge1:m^j\ge B_{\rm net}\},\quad C_0=128\,m\,B_{\rm net}^2
\end{equation}
(written $\nu$ in~\cite[\S5.4]{OAI26}, subscripted here to avoid collision with the bridge alias index $\nu\in\Z$ of
Section~\ref{sec:bridge}). At $h=100$ the exact integer values are
\begin{align*}
  W_bm-s_b&=2\,659\,138\,713\,000\,000,\qquad W_cm-s_c=6\,871\,402\,692\,000\,000,\\
  E_{\rm net}&=421\,070\,402\,185\,313\,436\,016\,911\,541\,181\,936\,724\,563\,007\,566\,900\,236\,040\,789\,184\,640\,064,\\
  B_{\rm net}&=421\,070\,402\,185\,313\,436\,016\,911\,541\,181\,936\,724\,564\,881\,374\,144\,872\,712\,056\,492\,640\,064,\\
  C_0&=22\,694\,436\,300\,352\,206\,084\,304\,065\,594\,397\,346\,702\,264\,594\,168\\
     &\phantom{{}={}}261\,607\,007\,508\,662\,043\,469\,612\,184\,775\,064\,463\,128\,060\,707\,693\\
     &\phantom{{}={}}386\,195\,338\,619\,369\,397\,054\,146\,227\,284\,214\,284\,288\,000\,000,
\end{align*}
where the three lines of $C_0$ concatenate to a $140$-digit integer ($C_0\approx2.2694\times10^{139}$,
$\log_2C_0=462.93\ldots$). Table~\ref{app:cert-closed} records the exact closed-form counts across representative
ground-set sizes $h$, showing the thresholds at which each family first achieves a rank saving ($h=39$ for the bit
network and $h=22$ for the complex network).

\begin{table}[htbp]
\centering\scriptsize
\caption{Exact closed-form counts ($m=h^3$) for the triple networks of~\cite[\S3]{OAI26}: roles $W$, edge-rank sum $s$,
signed saving $Wm-s$, and certified rational upper bound on $\ln m$. Only $h=100$ is used in the theorems.}
\label{app:cert-closed}\label{app:cert-ln}
\setlength{\tabcolsep}{3.5pt}
\begin{tabular}{@{}rlrrrl@{}}
\toprule
$h$ & network & $W$ & $s$ & $Wm-s$ & upper bound $\ln m$\\
\midrule
21 & bit & 3\,244\,397\,863\,700 & 30\,046\,370\,943\,598\,100 & $-$2\,327\,872\,400 & 9.133567313171\\
21 & complex & 6\,145\,204\,591\,400 & 56\,910\,739\,919\,072\,200 & $-$198\,116\,800 & 9.133567313171\\
22 & bit & 5\,628\,295\,349\,600 & 59\,930\,092\,117\,403\,200 & $-$3\,234\,862\,400 & 9.273127360075\\
22 & complex & 11\,249\,136\,760\,400 & 119\,780\,808\,120\,388\,800 & 104\,350\,400 & 9.273127360075\\
39 & bit & 4\,329\,455\,044\,049\,525 & 256\,818\,943\,756\,887\,996\,302 & 1\,085\,777\,173 & 10.990684938389\\
39 & complex & 16\,598\,761\,236\,611\,294 & 984\,621\,917\,049\,702\,208\,108 & 744\,843\,140\,678 & 10.990684938389\\
100 & bit & 177\,176\,569\,091\,445\,000\,000 & 177\,176\,569\,088\,785\,861\,287\,000\,000 & 2\,659\,138\,713\,000\,000 & 13.815510557965\\
100 & complex & 1\,873\,807\,244\,643\,542\,670\,000 & 1\,873\,807\,244\,636\,671\,267\,308\,000\,000 & 6\,871\,402\,692\,000\,000 & 13.815510557965\\
\bottomrule
\end{tabular}
\end{table}

\begin{proposition}[Network exponents at $h=100$]\label{prop:deficits}
For the integers $W_b,s_b,W_c,s_c$ of the triple networks of~\cite[\S3]{OAI26} at $h=100$ ($m=10^6$;
Definition~\ref{def:fixed-network-premises}),
\[
  \frac{s_b}{W_b}<m^{1-2^{-40}},\qquad \frac{s_c}{W_c}<m^{1-2^{-42}},
\]
so $\tau=1-2^{-40}$ and $\sigma=1-2^{-42}$ are admissible in Proposition~\ref{prop:prop14-param}. Moreover
$\nu_{\mathrm{net}}=11$, so $C_1=20$ and $C_1=14$ are both admissible. These are exact arithmetic properties of the
displayed integers; that the integers are the counts of working networks is Fact~\ref{fact:oai-networks}.
\end{proposition}

\begin{proof}
By Definition~\ref{def:fixed-network-premises} (\cite[\S3.6, (10)]{OAI26}) the relative deficits are
$\eta_b=339/22\,587\,335\,000\,000$ and $\eta_c=73/19\,906\,842\,167\,500$, with $s_b/W_b=m(1-\eta_b)$ and
$s_c/W_c=m(1-\eta_c)$. For $u>0$ and any $L\ge\ln m$,
\begin{equation}\label{eq:deficit-test}
  m^{1-u}=m\,e^{-u\ln m}>m(1-u\ln m)\ge m(1-Lu).
\end{equation}
Hence $s/W<m^{1-u}$ whenever $\eta\ge Lu$; with $u=2^{-e}$, it suffices to check $\eta\cdot2^e\ge L$. Take $L=14$:
since $\sum_{j=0}^{11}7^j/j!=41\,438\,755\,442/39\,916\,800>1000$, we have $e^7>10^3$ and hence $\ln 10^6<14$. The
exact integer comparisons
\[
  339\cdot2^{40}=372\,734\,441\,816\,064>316\,222\,690\,000\,000=14\cdot22\,587\,335\,000\,000,
\]
\[
  73\cdot2^{42}=321\,057\,395\,310\,592>278\,695\,790\,345\,000=14\cdot19\,906\,842\,167\,500
\]
give $\eta_b>14\cdot2^{-40}$ and $\eta_c>14\cdot2^{-42}$. For the guard degree, $10^{21}<W_c+m+1<2\cdot10^{21}$ and
$s_c<2\cdot10^{27}$, so $10^{60}<B_{\rm net}<512\cdot10^{63}+2\cdot10^{27}<10^{66}$; since $m^{10}=10^{60}$ and
$m^{11}=10^{66}$, we obtain $\nu_{\mathrm{net}}=11$.
\end{proof}

\begin{remark}[Sharper logarithm enclosure and sharpness]\label{rem:deficit-margins}
At $h=100$ the scaled deficits are $\eta_b2^{40}=16.50\ldots$ and $\eta_c2^{42}=16.12\ldots$, so the bound $\ln m<20$
used in~\cite[\S3.6]{OAI26} to certify $1-2^{-50}$ does not certify $1-2^{-40}$ or $1-2^{-42}$. Any certified upper
bound $L\le16.125$ suffices, and $L=14$ in Proposition~\ref{prop:deficits} is already elementary. A twelve-decimal
rational enclosure follows from $\ln10^6=36\operatorname{atanh}\frac13+12\operatorname{atanh}\frac19$: with
$S_N(z)=\sum_{j<N}z^{2j+1}/(2j+1)$ and $0<\operatorname{atanh}z-S_N(z)<z^{2N+1}/((2N+1)(1-z^2))$, taking $N=20$ yields
\[
  13.815510557964<\ln10^6<13.815510557965
\]
by finite rational comparisons (giving margins $\eta_b2^{40}-L=2.6864\ldots$ and $\eta_c2^{42}-L=2.3124\ldots$).
Moreover, the exponent $1-2^{-40}$ is false for the complex network at $h=100$: with $u=2^{-40}$ we have
$\eta_c2^{40}<5$ and $10<\ln m<20$, so $1-e^{-u\ln m}\ge u\ln m-\frac12u^2(\ln m)^2>u(10-200u)>5u>\eta_c$, whence
$s_c/W_c=m(1-\eta_c)>m^{1-2^{-40}}$ (and similarly for the bit network at $h=39$).
\end{remark}

\begin{lemma}[Elementary streams]\label{lem:transfer-streams}
Fix a radix $q_{\mathrm{rad}}\ge2$. The following operations have linear bit cost on a fixed number of tapes:
\begin{enumerate}[label=(\arabic*),leftmargin=*]
\item any fixed permutation of the pair of values at two specified radix-$q_{\mathrm{rad}}$ digit positions;
\item moving one specified radix-$q_{\mathrm{rad}}$ digit position to another position, preserving the order of the
  intervening positions;
\item splitting at a specified radix-$q_{\mathrm{rad}}$ digit, applying a fixed pointwise map to the aligned
  streams, and merging;
\item distributing complete rows cyclically among a fixed number $W$ of streams and, when the number of rows in each
  prefix is a multiple of $W$ (so that all streams have the same shape), merging them in a fixed order of the row
  classes;
\item exact fixed-width addition, subtraction, sign changes, and binary shifts by a supplied integer exponent $t$, when
  the necessary guard bits are present; a shift of $N$ words costs $O(N(w_{\rm in}+w_{\rm out})+\lg(2+|t|))$, linear in
  the input and output widths plus the length of the exponent's encoding, which is read once.
\end{enumerate}
Moreover, let a target field of range $[Q]$, $Q=q_{\mathrm{rad}}^b$, follow its control fields, each relevant control
having $O(b)$ digits. A cyclic shift of the target by an offset computable from the controls in $b^{O(1)}$ time costs
$O(V)$, and so does multiplication of the target coordinate by a fixed rational unit modulo $Q$ whose numerator and
denominator are prime to $q_{\mathrm{rad}}$.
\end{lemma}

\begin{proof}
\emph{Items~(1)--(4).} For~(1), split the array among $q_{\mathrm{rad}}^2$ tapes according to the two digit values;
within every stream the remaining address fields keep their order, so merging with the permuted pair of labels gives
the result. For~(2) and~(3), use $q_{\mathrm{rad}}$ streams and change the place in the nested merge at which the
selected digit is inserted. The number of streams is fixed, and complete spectator fields and trailing payloads are
copied without decoding. For~(4), send row $r$ to stream $r\bmod W$; each stream keeps the row order, and a merge reads
the streams in the prescribed order of row classes, copying every row a constant number of times.

\emph{Overheads and item~(5).} Binary counters are stored least significant bit first. An increment, including the
return of the head over the changed prefix, has constant amortized cost, and wraparound and reset costs are charged to
the preceding traversal. Initial shape preparation takes at most $C(\log_2(2V))^{J}$ steps for fixed $C,J$, which is
at most $C'V$ for all $V\ge1$. This accounting covers counters, head returns and clearing the regions visited by a
call. Fixed-width coefficient arithmetic uses ordinary carry scans, and a bit reversal inside a coefficient, when
needed, is a charged linear pass. For a shift by a supplied exponent $t$, partition the source and destination words
into the retained interval and the intervals of known zeros or redundant sign bits; then copy the retained interval and
fill the rest in one traversal. Exact right shifts discard only known zero low bits, and local counter preparation
costs $O(\log(2w))$ per $w$-bit word, absorbed by its scan. Thus a shift by $t$ is executed in one traversal rather
than $t$ one-bit shifts, as~\cite[\S5.4, p.~36]{OAI26} also notes. This gives~(5).

\emph{Controlled shifts and units.} Fix the preceding coordinates. The target fiber consists of $Q$ consecutive suffix
blocks, and a rotation splits it at the computed cut and concatenates the two intervals in the opposite order. The
offset computation is polynomial in $b$, whereas the fiber contains at least $q_{\mathrm{rad}}^b$ bits, so the fiber
scan absorbs the offset work, uniformly in $b$. Multiplication by a fixed positive integer $a$ prime to $Q$ splits the
target coordinates into the at most $a$ intervals on which $x\mapsto ax-kQ$ lies in $[Q]$; on each interval the outputs
increase and lie in one residue class modulo $a$, and these classes are distinct since $\gcd(a,Q)=1$. Merging the fixed
number of streams in output order gives the product, and the inverse uses the reversed schedule. A negative multiplier
adds the negation $y\mapsto-y\bmod Q$, which fixes the block at $y=0$ and reverses the order of the other $Q-1$ blocks
while preserving each block's internal order: push the blocks onto a temporary tape and reverse each popped block once
more, as in~\cite[proof of Lemma~4, \S2.1]{OAI26}. A rational unit is an integer multiplication composed with an
inverse one. We use the fixed-tape details of~\cite[Lemmas~3--4, \S2.1]{OAI26} as summarized in
Fact~\ref{fact:oai-streams} (applied in~\cite[proof of Lemma~10, proof of Prop.~11, \S5.2, and \S5.3 (25)]{OAI26}).
\end{proof}

The next lemma justifies the address-dependent Walsh diagonals $S_U$ and $\CZ$ of Section~\ref{sec:walsh} and the
twiddle-index preparation of Section~\ref{sec:direct-fft}.

\begin{lemma}[Long-record address-dependent scans]\label{lem:transfer-long-record}
Fix a recordwise address routine and payload routine and constants $C_0,C_1,C_2,J$. Suppose that, on an array of
$M\ge1$ records, the address preparation of every record takes at most $C_1q^J$ steps, the payload processing of a
record of $\rho$ bits takes at most $C_2\rho$ steps, and a one-time setup takes at most $C_0q^J$ steps. Then the total
time is at most $C_0q^J+M(C_1q^J+C_2\rho_{\max})+O(V)$; in particular, if every record has at least $q^J$ bits (which
holds for records of $\Theta(rq)$ bits once $r\ge q^J$, as in every family with $r/q^j\to\infty$ for each fixed $j$),
the total time is $O(V)$ with a constant depending only on $C_0,C_1,C_2$. Consequently, an address-dependent phase in
$\{1,-1,i,-i\}$ may be computed once per record and applied to all its coefficients in $O(V)$ time, and selecting a
twiddle from a table of $O(rq)$ bits by one scan per record has $O(V)$ overhead (with general coefficient
multiplications charged separately).
\end{lemma}

\begin{proof}
Let the records have $\rho_1,\dots,\rho_M$ bits, so that $V=\rho_1+\dots+\rho_M\le M\rho_{\max}$. Summing over the
records, setup, address preparation and payload processing take at most $C_0q^J+MC_1q^J+C_2V$ steps, and traversing the
array on a fixed number of tapes adds $O(V)$. When every record has at least $q^J$ bits, $q^J\le V$ and $Mq^J\le V$,
giving $O(V)$. For records of $\Theta(rq)$ bits, maintaining and decoding the record address, scanning a polynomial-size
mask or position list, and evaluating a fixed polynomial-time address routine cost at most $Cq^J$ per record, so the
total is $O(M(q^J+rq))=O(Mrq)=O(V)$, and a table scan and rewind cost $O(rq)$ per record. The same long-record
amortization appears in~\cite[\S5, p.~26 and \S5.3, pp.~33--34]{OAI26}.
\end{proof}

\subsection{Rational shears and power-width interchanges}\label{sec:transfer-power}

For chunks of $m^k$ digits, the bit network, with the shears of Lemma~\ref{lem:transfer-shear} as frame changes,
reduces one interchange to $s_b$ interchanges of chunks $m$ times shorter, each on $1/W_b$ of the volume, plus linear
work (Proposition~\ref{prop:transfer11}).

\begin{lemma}[Lower triangular factorization and shear]\label{lem:transfer-shear}
Every rational $m\times m$ matrix $A$ of rank $a$ has a factorization $A=E_1\Pi E_2$, where $E_1,E_2$ are invertible
lower triangular rational matrices and $\Pi$ is a partial permutation matrix with exactly $a$ ones. For a fixed finite
collection of such matrices there is a fixed odd prime $q_{\mathrm{rad}}$ such that, on two ordered groups
$H=(H_1,\dots,H_m)$ and $D=(D_1,\dots,D_m)$ of fields of range $[q_{\mathrm{rad}}^b]$ with every $H_i$ before every
$D_j$ and arbitrary spectator fields, the shear $(H,D)\mapsto(H+AD,D)$ uses exactly $a$ interchanges of pairs of
$b$-digit fields, each between an $H$ field and a $D$ field, and $O(V)$ further work.
\end{lemma}

\begin{proof}
\emph{Factorization.} If the active matrix is nonzero, choose its topmost nonzero row and the rightmost nonzero entry
of that row as the pivot. Right multiplication by lower triangular elementary matrices clears entries to the left of
the pivot, and left multiplication by lower triangular elementary matrices clears entries below it. Scale the pivot to
one and delete its row and column from the active index sets, retaining their original orders. Subsequent operations
remain lower triangular and preserve earlier pivots, so after $a$ steps the result is a partial permutation matrix with
$a$ ones; inverting the accumulated triangular matrices gives $A=E_1\Pi E_2$ (\cite[Lemma~9, \S4.1]{OAI26}).

\emph{The shear.} Fix the factorizations for the finite collection once and for all, and choose an odd prime
$q_{\mathrm{rad}}$ avoiding all denominators of the factors and their inverses and all numerators of their diagonal
entries by finite rational arithmetic and trial division (\cite[proof of Lemma~10, \S4.1]{OAI26}). Every identity then
specializes modulo $q_{\mathrm{rad}}^b$ with invertible triangular factors. A lower triangular map costs $O(V)$ by
updating coordinates in descending physical order using only earlier, still-unmodified coordinates
(Lemma~\ref{lem:transfer-streams}, since $m$ is fixed). Transform $D$ by $E_2$ and $H$ by $E_1^{-1}$, add $\Pi D$ to
$H$, and undo both triangular changes. Each one of $\Pi$ at $(i,j)$ is executed by
\[
  D_j\leftarrow D_j+H_i,\qquad(H_i,D_j)\leftarrow(D_j,H_i),\qquad D_j\leftarrow H_i-D_j ,
\]
which maps $(h,d)\mapsto(h+d,d)$ using one interchange and two linear updates to the later field
(\cite[proof of Lemma~10, \S4.1]{OAI26}).
\end{proof}

\begin{proposition}[Parameterized power-width interchange]\label{prop:transfer11}
Take the bit network of Definition~\ref{def:fixed-network-premises}, and choose the prime $q_{\mathrm{rad}}$ of
Lemma~\ref{lem:transfer-shear} for the edge matrices $A_e$, also avoiding the denominators of the assigned terminal and
gate matrices. For every fixed rational $\tau\in(0,1)$ with $s_b/W_b<m^\tau$, two disjoint contiguous chunks of $e=m^k$
radix-$q_{\mathrm{rad}}$ digits can be interchanged in $O(Ve^\tau)$ time if a row field preceding both chunks has
length divisible by $W_b^k$. All other fields are preserved, and the machine, tape count and time constant are uniform
in all field lengths and data.
\end{proposition}

\begin{proof}
\emph{Construction and frame invariant.} Write $W=W_b$ and $s=s_b$. The case $k=0$ is Lemma~\ref{lem:transfer-streams}(1).
For $k>0$, write each chunk as $m$ fields of $b=e/m$ digits, and split whole rows cyclically among $W$ role streams:
writing the old row index as $Wg+(w-1)$, role $w$ has row index $g$, volume $V/W$, and row count divisible by
$W^{k-1}$. For a rational matrix $M$ let $\Phi_M$ move $(H,D)$ to $(H+MD,D)$, so $\Phi_{M'}\Phi_M=\Phi_{M'+M}$ and
$\Phi_M^{-1}=\Phi_{-M}$. Apply $\Phi_{A_e}$ on each circuit edge by Lemma~\ref{lem:transfer-shear} with recursive pivot
interchanges, and apply each circuit gate pointwise to aligned role streams. Defining the logical array at a vertex
with frame $\Phi_M$ as $\Phi_M^{-1}$ of the physical array, edges preserve the logical array and common address
permutations commute with pointwise gates, so input role $w$ reaches output role $\rho_b(w)$ with physical operator
$\Phi_{M_{\mathrm{out}(\rho_b(w))}}\Phi_{-M_{\mathrm{in}(w)}}=\Phi_{I_m}$. Merging output role $\rho_b(w)$ back into row
class $w$ applies $H\leftarrow H+D$ on every original row, including auxiliary rows with arbitrary initial contents
(\cite[proof of Prop.~11, \S4.2]{OAI26}). Then $D\leftarrow D-H$, $H\leftarrow H+D$, $D\leftarrow H-D$ executes
$(H,D)\mapsto(D,H)$ with one circuit call and two $O(V)$ updates (\cite[proof of Prop.~11, pp.~22--23]{OAI26}).

\emph{Recurrence, fixed tapes, and descriptors.} There are $s$ recursive calls on volume $V/W$ plus $O(V)$ work, so
$F_0=O(1)$ and $F_k\le(s/W)F_{k-1}+O(1)$; since $a=s/W<m^\tau$, $\sum_{j=0}^ka^j=O(m^{k\tau})$ and $F_k=O(e^\tau)$
(\cite[proof of Prop.~11, (12)]{OAI26}). Depth-first scheduling uses $W$ role tapes, fixed I/O and work tapes, a
parking-stack tape and a descriptor-stack tape. Inactive streams are pushed before a child call, parent role areas are
erased while traversed, and popped streams are written from the right end of their known destination intervals to
restore direction without scanning ancestor parked data, adding $O(V)$ at the parent (\cite[\S4.2, p.~23]{OAI26}). At
depth $j$ from root width $e_0=m^{k_0}$, inactive digits remain spectators, so $V_j\ge q_{\mathrm{rad}}^{2e_0}$ and
$\log V_0=O(\log(2V_j))$, absorbing the $\mathrm{poly}(\log(2V_j))$ child-descriptor work into $O(V_j)$
(\cite[\S4.2, pp.~23--24]{OAI26}).
\end{proof}

\subsection{Arbitrary-width interchanges}\label{sec:transfer-arbitrary}

\begin{lemma}[Parameterized arbitrary-width interchange]\label{lem:transfer12}
Under the hypotheses of Proposition~\ref{prop:transfer11}, for the same fixed $\tau$, a bit array on
$[A]\times[2^u]\times[G]\times[2^u]\times[B]$, $u\ge1$, can be transformed by $(x,h,g,d,z)\mapsto(x,d,g,h,z)$ in
$O(Vu^\tau)$ time, uniformly in all positive $A,G,B$ and all data (including shape processing, padding, unpadding and
cleanup, with $B=1$ permitted). Two fields whose widths differ by one bit are interchanged across the intervening field
by one such interchange and one single-bit move.
\end{lemma}

\begin{proof}
\emph{Radix-$q_{\mathrm{rad}}$ chunks.} First let the chunks have range $[q_{\mathrm{rad}}^e]$. Put
$k=\lceil\log_me\rceil$ and $k_{\rm hi}=\min\{j\ge0:q_{\mathrm{rad}}^{2j}\ge W_b^k\}$. For large $e$,
$1\le k_{\rm hi}<e$ with $k_{\rm hi}=O(\log(2e))$ (bounded widths use individual digit swaps). Split each chunk into
its top $k_{\rm hi}$ digits and the rest, $h=(h_+,h_-)$ and $d=(d_+,d_-)$. Interchange $h_+$ and $d_+$ digit by digit
and move $h_+$ immediately after $d_+$, transforming $(h_+,h_-,g,d_+,d_-)\mapsto(d_+,h_+,h_-,g,d_-)$ in $O(Vk_{\rm hi})$
steps so that $(d_+,h_+)$ forms a row field of range $R_{\rm row}=q_{\mathrm{rad}}^{2k_{\rm hi}}\ge W_b^k$ before both
remaining chunks. Pad this row range in each prefix to $R'_{\rm row}=W_b^k\lceil R_{\rm row}/W_b^k\rceil<2R_{\rm row}$.
Expand $e-k_{\rm hi}=\sum_{j=0}^ka_jm^j$ ($0\le a_j<m$) and interchange corresponding power-of-$m$ pieces by
Proposition~\ref{prop:transfer11} in cost $O(V)\sum_ja_jm^{j\tau}=O(Ve^\tau)$. Delete the padded rows and reverse only
the joining moves to reach $(d_+,d_-,g,h_+,h_-)$, with descriptor and movement overhead $O(V\log(2e))=O(Ve^\tau)$.
These steps follow~\cite[proof of Lemma~12, \S4.3]{OAI26}; none uses the specific value $1-2^{-50}$ fixed
in~\cite[\S4.2, p.~24]{OAI26}.

\emph{Binary chunks.} Choose the least $e$ with $q_{\mathrm{rad}}^e\ge2^u$ and pad both ranges numerically from $[2^u]$
to $[q_{\mathrm{rad}}^e]$ (increasing volume by less than $q_{\mathrm{rad}}^2$ without per-address radix conversion, as
numerical and lexicographic orders coincide). After the radix-$q_{\mathrm{rad}}$ interchange, scan away entries with
either coordinate $\ge2^u$; since $e=O(u)$, the time is $O(Vu^\tau)$ (\cite[proof of Lemma~12, p.~25]{OAI26}). Widths
differing by one bit use one equal-width interchange and one single-bit move as in~\cite[\S4.3, p.~25]{OAI26}.
\end{proof}

\subsection{Packed selected-bit addition}\label{sec:transfer-packed}

\begin{lemma}[Parameterized packed selected-bit addition]\label{lem:transfer13}
Fix $\tau$ as in Lemma~\ref{lem:transfer12}. Let $R_{\mathrm{rec}}\ge1$, $K\ge6$, $f\ge1$, $L=fK$ and $0\le\rho<K$, let
$N_0,\dots,N_3$ be positive integers, and let $(\chi_1,\chi_2,\chi_3)$ be a permutation of the chunk names $(x,y,z)$.
Consider $R_{\mathrm{rec}}$-bit records on the complete rectangle
$[N_0]\times[2^L]_{\chi_1}\times[N_1]\times[2^L]_{\chi_2}\times[N_2]\times[2^L]_{\chi_3}\times[N_3]$, and put
$M=2^{3L}\prod_{j=0}^3N_j$, $V=MR_{\mathrm{rec}}$ and $A_{\mathrm{rect}}=\lceil\log_2(2V)\rceil$. One fixed-tape
procedure performs $y_{\rho+iK}\leftarrow y_{\rho+iK}\oplus x_{\rho+iK}$ for $0\le i<f$, preserves every other address
bit and every payload, and restores the entire auxiliary chunk $z$ for every initial address value. Its time is
\begin{equation}\label{eq:transfer-packed}
  O\bigl(V(L^\tau+1)+MA_{\mathrm{rect}}^3+\delta_{\mathrm{bad}}MA_{\mathrm{rect}}(R_{\mathrm{rec}}+A_{\mathrm{rect}})\bigr),
  \qquad \delta_{\mathrm{bad}}=\min\{1,80(f-1)2^{-K}\}.
\end{equation}
In particular, if $A_{\mathrm{rect}}^3\le R_{\mathrm{rec}}$ and
$80(f-1)2^{-K}A_{\mathrm{rect}}(1+A_{\mathrm{rect}}/R_{\mathrm{rec}})\le1$, the bound is $O(V(L^\tau+1))$ with an
absolute implied constant; both hold for all sufficiently large $q$, uniformly in shapes and payloads, whenever
$A_{\mathrm{rect}}\le C_Aq$, $f\le d\le q$, $R_{\mathrm{rec}}\ge q^4$ and $K\ge3\log_2q+\log_2(160\,C_A)$.
\end{lemma}

\begin{proof}
\emph{Scalar program and packing.} For $c\in\{0,1\}$ and integers $u,w$, execute in chronological order (congruences
modulo~$2$, using current values):
\[
\begin{array}{llll}
u\leftarrow u-c[w\equiv0], & w\leftarrow w-c[u\equiv0], & u\leftarrow u+c[w\equiv0], & w\leftarrow w+c[u\equiv0],\\
w\leftarrow w+c[u\equiv1], & u\leftarrow u+c[w\equiv1], & w\leftarrow w-c[u\equiv1], & u\leftarrow u-c[w\equiv1].
\end{array}
\]
For $c=0$ this is the identity; for $c=1$ the two halves induce $(00\ 01\ 10)$ and $(01\ 10\ 11)$ on each square
$\{2u_0,2u_0+1\}\times\{2w_0,2w_0+1\}$, composing to $(u',w')=(u+c(1-2(u\bmod2)),\,w)$ with displacement at most four
(\cite[proof of Lemma~13, (20)--(21)]{OAI26}). Omitting the top selected bit $j_{f-1}$, write the $K$-bit segments at
$j_i=\rho+iK$ ($i<f-1$) in $y,z$ as $u_i=2g_i+a_i$, $w_i=2h_i+b_i$. Each program line across all $i<f-1$ is a single
cyclic shift modulo $2^L$ (e.g.\ $y\leftarrow y+\epsilon\sum_{i=0}^{f-2}2^{j_i}[x_{j_i}=1][z_{j_i}=b]\bmod2^L$), hence
a bijection of the rectangle. Let $B$ be the set of addresses where some guard $g_i$ or $h_i$ fails
$10\le g\le2^{K-1}-11$. Off $B$, every segment starts in $[20,2^K-21]$ and stays within $[0,2^K-1]$ without inter-segment
carries, so the eight-shift composite $S_0$ agrees with the ideal XOR map $T_0$ on $i<f-1$. By the union bound over
$2(f-1)$ guards with $20$ excluded values out of $2^{K-1}$ each, $|B|\le\delta_{\mathrm{bad}}M$; since $T_0$ preserves
guards, $T_0S_0^{-1}$ permutes $B$ and fixes $B^c$ (\cite[proof of Lemma~13, (22)--(24)]{OAI26}).

\emph{Movement and deterministic repair.} Moving each shift's target to the rightmost chunk position by
Lemma~\ref{lem:transfer12}, applying the controlled rotation of Lemma~\ref{lem:transfer-streams}, and undoing the swap
uses at most $16$ interchanges and $8$ rotations at cost $O(V(L^\tau+1))$ (\cite[proof of Lemma~13, pp.~30--31]{OAI26}).
Next scan with a record-rank counter, test membership in $B$ in $O(A_{\mathrm{rect}}^3)$ steps per record, extract
records $q_{\rm addr}\in B$ with destination rank of $T_0(S_0^{-1}(q_{\rm addr}))$, mark their holes, radix-sort the
extracted bank in $O(|B|A_{\mathrm{rect}}(R_{\mathrm{rec}}+A_{\mathrm{rect}}))$ time, and reinsert in one forward scan;
finally apply the omitted two-bit XOR at $j_{f-1}$ by Lemma~\ref{lem:transfer-streams}(1) in $O(V)$ time
(\cite[proof of Lemma~13, (19), p.~31]{OAI26}). Neither $B$ nor the repair depends on $\tau$.
\end{proof}

\subsection{Complex frames and the simultaneous butterfly layer}\label{sec:transfer-frames}

Put $H_{\mathrm u}=\sqrt2H_0$, $S=\diag(1,i)$, $Z=\diag(1,-1)$ and $C=H_{\mathrm u}SH_{\mathrm u}=aI+bX$ with
$a=\frac{1+i}2$, $b=\frac{1-i}2$ and $X=\bigl(\begin{smallmatrix}0&1\\1&0\end{smallmatrix}\bigr)$. Direct multiplication
gives $H_0=bSCS$, $C^{-1}=-iZCZ$ and $b^2=-i/2$ (\cite[\S5.1, (13)--(14) and \S5.4, p.~36]{OAI26}).

\begin{lemma}[Exact recursive complex contract]\label{lem:transfer-phase}
For the complex network of Definition~\ref{def:fixed-network-premises}, on a group of $e=mf$ consecutive $K$-bit chunks
with one common selected offset, $C^{\otimes e}$ is obtained on every role, exactly over $\Z[i,1/2]$, by
\begin{enumerate}[label=(\arabic*),leftmargin=*]
\item exactly $s_c$ forward child calls $C^{\otimes f}$, each on one slot of $f$ consecutive chunks;
\item a fixed number of packed binary row additions and their inverses;
\item a fixed number of coefficientwise Gaussian dyadic scalar gates, address signs and fourth-root phases.
\end{enumerate}
No extra child call is needed for an inverse directional factor or for an endpoint correction.
\end{lemma}

\begin{proof}
For a binary subspace label $U$ put $q_U(x)=\operatorname{wt}(P_Ux)\bmod4$ and
$C_U=H_{\mathrm u}^{\otimes m}\diag_x(i^{q_U(x)})H_{\mathrm u}^{\otimes m}$. Orthogonal binary vectors meet in an even
number of coordinates, so weight is additive modulo~$4$ on orthogonal sums. On a nested edge with orthonormal residual
basis $\mathcal B$, $q_V(x)-q_U(x)=\sum_{v\in\mathcal B}\operatorname{wt}(v)[v\cdot x]\bmod4$ (negated on a decreasing
edge), with each $\operatorname{wt}(v)$ odd. Since $i^{[v\cdot x]}=a+b(-1)^{v\cdot x}$ and conjugating
$\diag_x((-1)^{v\cdot x})$ by $H_{\mathrm u}^{\otimes m}$ yields $(X_vf)(x)=f(x+v)$, each $v\in\mathcal B$ contributes
$aI+bX_v$ or its inverse (\cite[proof of Prop.~8, \S3.5]{OAI26}).

Extending $\mathcal B$ to $M\in\mathrm{GL}(m,\mathbb F_2)$ with first columns $\mathcal B$, the move permutation $P_M$
satisfies $P_MX_vP_M^{-1}=X_{Mv}$, realizing the edge on $f$ columns by $P_M^{-1}$, one coordinate kernel per
$v\in\mathcal B$, and $P_M$, where inverse factors use $(C^{-1})^{\otimes f}=(-i)^fZ^{\otimes f}C^{\otimes f}Z^{\otimes f}$
without extra child calls (\cite[\S5.1, (15)--(17)]{OAI26}). Each fixed $M^{\pm1}$ factors into at most $m(m-1)+3m$ row
additions, each executed across two $f$-chunk slots by Lemma~\ref{lem:transfer13} using a third slot ($m\ge3$) as
restored scratch (\cite[\S5.1, p.~27 and \S5.2, pp.~31--32]{OAI26}). Commutation of common address-linear frames with
pointwise role gates, together with the source/sink signs $Z_{a,f}$, scalar $i^{27f}$ and bank restoration, yields
$C^{\otimes e}$ on every role from $s_c$ child calls (\cite[proof of Prop.~8, \S3.5 and \S5.1, pp.~27--28]{OAI26}).
Finally,
\begin{equation}\label{eq:oai-h0-transfer}
  H_0^{\otimes D}=b^D\mathcal S_DC^{\otimes D}\mathcal S_D,\qquad
  (\mathcal S_Dz)(x)=i^{\sum_{j=1}^Dx_j}z(x),\qquad b^D=(-i/2)^{\lfloor D/2\rfloor}b^{D\bmod2},
\end{equation}
as in~\cite[\S5.1, (18)]{OAI26}, with $\mathcal S_D$ executed in $O(V)$ time by Lemma~\ref{lem:transfer-long-record}.
\end{proof}

\begin{proposition}[Simultaneous butterfly layer with general parameters]\label{prop:prop14-param}
Let $\nu_{\mathrm{net}}$ and $C_0$ be as in~\eqref{eq:walsh-guard-constants}. Fix rationals $0<\tau<1$ and $0\le\sigma<1$ with
\[
  s_b/W_b<m^{\tau},\qquad s_c/W_c<m^{\sigma},
\]
and rational parameters $0<\beta<1$, $c>0$, $0<\eps<1$, $\lambda$, $\lambda'$ with
\begin{align}
  \max\{\tau,\sigma\}<\lambda<1,\qquad \tau(1+c/\beta)<\lambda,\label{eq:walsh-layer-a}\\
  \max\{\lambda,\ \sigma+\beta(1-\sigma)\}<\lambda'<1.\label{eq:walsh-layer-b}
\end{align}
Fix an integer $C_1\ge\nu_{\mathrm{net}}+3$ with $\eps C_1<1$, and positive constants $a_d\le b_d$, $a_r\le b_r$, $C_M$,
$C_{\rm desc}$ and $C_w>1$. Then there are a cutoff $q_{\min}$, one fixed finite-alphabet multitape procedure and a
uniform time constant with the following property. Suppose that
\begin{equation}\label{eq:walsh-layer-bands}
  q\ge q_{\min},\qquad a_dq^{\eps}\le d\le b_dq^{\eps},\qquad 2^{a_rq/d}\le r\le 2^{b_rq/d},\qquad K=\lfloor d^c\rfloor,
\end{equation}
and let $1\le D\le d$, $A,S\ge1$ and $0\le\rho<K$ be integers with
\[
  M=AS\,2^{KD},\qquad \lg(2M)\le C_Mq,\qquad q+2\le w\le C_wq.
\]
The input is a disk-grid array in lexicographic order on $[A]\times[2^K]^D\times[S]$. Each address carries a record in
$\Rr_r$, stored in coefficient order, and each component is $2^{-q}$ times a signed $w$-bit two's-complement integer,
most significant bit first. The supplied header is
$\Gamma(q)\Gamma(d)\Gamma(r)\Gamma(D)\Gamma(A)\Gamma(S)\Gamma(K)\Gamma(\rho+1)\Gamma(w)$, of length at most
$C_{\rm desc}q$, with $\Gamma$ as in~\eqref{eq:oai-gamma}. Let $H$ apply $H_0$ to bit~$\rho$, counted from the least
significant end, of each of the $D$ chunks. The procedure returns $Q_q(Hz)$, in the original record and component
format, in $O(Vd^{\lambda'})$ steps, where $V=2Mrw$. It evaluates the entire layer exactly before this final
truncation; the output is a disk array and $\|Q_q(Hz)-Hz\|_\infty<\sqrt2\cdot2^{-q}$. The recursion stops when
$e<d^\beta$ (with $e=1$ always a leaf). With $\Delta=\lceil C_0d^{C_1}\rceil$, the procedure uses at most $q+\Delta$
fractional and $\Delta+3$ integer and sign bits per real component internally; these widths are $O(q)$. The time
includes padding, deterministic repair, role-stream scheduling, unpadding and workspace cleanup.
\end{proposition}

\begin{remark}[Where the parameters enter]\label{rem:prop14-ledger}\label{rem:walsh-delta}\label{rem:transfer-guard}
As printed, \cite[Prop.~14]{OAI26} fixes $\beta=\frac12$, $\tau=\sigma=1-2^{-50}$ and $C_1=20$. In the proof below,
following~\cite[\S\S5.1--5.4]{OAI26}, the parameters enter only through the hypotheses of
Proposition~\ref{prop:prop14-param}. The parameter $\beta$ enters the leaf test $e<d^\beta$ (evaluated for
$\beta=\beta_1/\beta_2$ as the integer comparison $e^{\beta_2}<d^{\beta_1}$) and the bound $K\le e^{c/\beta}$. The
exponents $\tau$ and $\sigma$ enter the interchange costs of Lemmas~\ref{lem:transfer12}--\ref{lem:transfer13} and the
time unrolling via $s_c/W_c<m^\sigma$, respectively. Finally, $C_1\ge\nu_{\mathrm{net}}+3$ enters the actual guard
allocation $\Delta=\lceil C_0d^{C_1}\rceil$ together with $\eps C_1<1$, which ensures $\Delta=o(q)$ (already noted
in~\cite[\S5.4, p.~35]{OAI26} as requiring only exponent $14=\nu_{\mathrm{net}}+3$).
\end{remark}

\subsection{Proof of Proposition~\ref{prop:prop14-param}}\label{sec:transfer-layer}

\emph{Individual kernels.} Applying $C(u,v)=\bigl(\frac{u+v+i(u-v)}2,\frac{u+v-i(u-v)}2\bigr)$ at one selected bit splits
on that bit, scans aligned coefficient pairs, and merges in $O(V)$ time at any $O(q)$ guarded width, so $e$ individual
kernels cost $O(Ve)$ (\cite[\S5.3, (25)]{OAI26}).

\emph{Rows, pieces, and scheduling.} Put $d_0=\lceil\log_2W_c\rceil\lceil\log_m(2d)\rceil$ and $k_0=\lceil\log_md\rceil$.
If $D\le d_0$, use individual kernels. Otherwise apply individual kernels to the first $d_0$ selected bits; their
complete consecutive chunks form a row field of range $R_{\mathrm{row}}=2^{d_0K}\ge W_c^{k_0}$, padded in each prefix to
the next multiple of $W_c^{k_0}$ (at most doubling volume), and the remaining $D-d_0$ chunks split into at most
$(m-1)(1+\lfloor\log_md\rfloor)$ consecutive power-of-$m$ pieces. At each internal node, cyclic row splitting gives $W_c$
role streams of volume $V/W_c$ with row count divisible by $W_c^{k_0-j}$ at depth $j\le k_0$, supporting recursion down
to one chunk regardless of $\beta$ and preserving the three-slot rectangles of Lemma~\ref{lem:transfer13}
(\cite[\S5.3, pp.~33--34]{OAI26}). Depth-first parking as in Proposition~\ref{prop:transfer11} adds $O(V)$ overhead per
node across the fixed number $s_c$ of children, and completed invocations commute with the row-bit kernels and leave
padded zero rows zero for removal (\cite[\S5.3, pp.~33--34]{OAI26}).

\emph{Arithmetic depth and guard.} Let $E_*$ be the constant number of elementary coefficient operations (additions,
subtractions, halvings, signs and multiplications by~$i$) performed outside child calls at one internal node by the
fixed finite scalar circuit and its fourth-root/sign corrections. Then the dependency depth satisfies
$A(e)\le s_cA(e/m)+E_*$ at internal nodes and $A(e)\le8e$ at leaves (\cite[\S5.4, p.~35]{OAI26}). Over
$j_{\max}\le\lceil\log_md\rceil$ levels with $1<s_c<m^5$, we have $s_c^{j_{\max}}\le s_cd^5$ and
$A(e)\le C_{\rm dep}d^6$ where $C_{\rm dep}=s_c(8+E_*/(s_c-1))$. Across at most $md$ pieces (for large $d$), at most $d$
preprocessing kernels, the two outer diagonals $\mathcal S_D$ and the scalar $b^D$ (whose shift by $t=\lfloor D/2\rfloor$
is charged as $t$ halvings), total denominator and magnitude growth obeys (\cite[\S5.4, p.~36]{OAI26})
\begin{equation}\label{eq:walsh-guard-growth}
  G(d)\le mC_{\rm dep}d^7+8d+t+9\le(mC_{\rm dep}+18)\,d^7\le C_0d^{C_1}\qquad(d\ge d_*),
\end{equation}
with $d_*=\bigl\lceil((mC_{\rm dep}+18)/C_0)^{1/(C_1-7)}\bigr\rceil$ finite since $C_1\ge\nu_{\mathrm{net}}+3=14>7$; the condition
$d\ge d_*$ holds via $d\ge a_dq^\eps$ once $q\ge q_{\min}$. Starting from disk inputs on $2^{-q}\Z[i]$, values at depth
$j$ lie in $2^{-(q+j)}\Z[i]$ with modulus at most $2^j$, so $\Delta=\lceil C_0d^{C_1}\rceil$ ($q+\Delta$ fractional and
$\Delta+3$ integer/sign bits) represents every intermediate value exactly, and temporary bit-level interchanges cause no
coefficient rounding (\cite[\S5.4, p.~36]{OAI26}). Since $d\le b_dq^\eps$ and $\eps C_1<1$, $\Delta=o(q)$, so
$\Delta\le q$ for $q\ge q_{\min}$ and guarded records have width $R_g=2r(q+2\Delta+3)\in[2rq,12rq]$.

\emph{Uniform packed-call cost and time recurrence.} Every packed child has
$A_*=\lceil\log_2(2M_{\rm child}R_g)\rceil\le C_Aq$ with $C_A=C_M+b_r+7$, while $K\ge\frac12a_d^cq^{\eps c}$ and
$r\ge2^{(a_r/b_d)q^{1-\eps}}$, so both extra normalized terms of~\eqref{eq:transfer-packed} tend to zero uniformly across
all children and descriptors/masks are absorbed by Lemma~\ref{lem:transfer-long-record}
(\cite[\S5.3, pp.~33--34 and proof of Prop.~14, \S5.4, pp.~37--38]{OAI26}). With $a_c=s_c/W_c$ and $F(e)$ the time
divided by current logical volume,
\begin{equation}\label{eq:walsh-layer-recurrence}
  F(e)\le a_cF(e/m)+O((eK)^{\tau}+1),\qquad F(e)=O(e)\quad(e<d^\beta).
\end{equation}
For $e\ge d^\beta$ we have $K\le d^c\le e^{c/\beta}$, so~\eqref{eq:walsh-layer-a} bounds the node overhead by
$O(e^\lambda)$ and internal nodes sum to $O(e^\lambda)$ since $a_c<m^\sigma<m^\lambda$, while leaves of size
$u=e/m^j<d^\beta$ contribute $O(a_c^ju)\le O(e^\sigma u^{1-\sigma})\le O(e^\sigma d^{\beta(1-\sigma)})$. Summing over
the base-$m$ pieces and preprocessing yields
\begin{equation}\label{eq:walsh-layer-unrolled}
  O\bigl(\log(2d)\,[d^{\lambda}+d^{\sigma+\beta(1-\sigma)}+1]\bigr)=O(d^{\lambda'})
\end{equation}
by~\eqref{eq:walsh-layer-b} (\cite[\S5.3, p.~34]{OAI26}).

\emph{Order of operations and output.} Re-encode into the guarded format (sign-extending or trimming redundant sign bits
to $\Delta+3$ integer/sign bits and appending $\Delta$ fractional zero bits, giving exact width $q+2\Delta+3$), apply
$\mathcal S_D$ in the original layout, execute the row-bit kernels, row padding and power-of-$m$ pieces for
$C^{\otimes D}$, unpad, apply $\mathcal S_D$ again and $b^D$ as in~\eqref{eq:oai-h0-transfer}, and apply $Q_q$ once.
The chronological order matters: for $D=1$, placing both $\mathcal S_1=S$ factors after $C$ would give
$bS^2C(1,0)^{\mathsf T}=(\frac12,\frac i2)^{\mathsf T}\ne H_0(1,0)^{\mathsf T}$. The final truncation acts on the
absolute numerator and then restores the sign, implementing componentwise truncation toward zero $Q_q$ in two's
complement; because each row of $H=H_0^{\otimes D}$ has absolute sum~$1$, $H$ preserves the unit disk, and $Q_q$
preserves the disk with error $<\sqrt2\cdot2^{-q}$ in width $w\ge q+2$ (\cite[proof of Prop.~14, p.~38]{OAI26}). Taking
$q_{\min}$ above the finitely many thresholds proves Proposition~\ref{prop:prop14-param}.\qed

\begin{remark}[Explicit gate-count threshold]\label{rem:transfer-guard-explicit}
The proof above uses only finiteness of the scalar circuit. Counting at most $12W_c$ scalar gates of at most $2W_c^2$
operations each, $4s_c$ inverse-child operations and $4W_c+4$ endpoint operations as in the proof text
of~\cite[\S5.4, pp.~35--36]{OAI26} gives $E_*\le24W_c^3+4s_c+4W_c+4<E_{\rm net}$, whence
$A(e)\le8B_{\rm net}^2d^{\nu_{\mathrm{net}}+1}$ and $G(d)<128mB_{\rm net}^2d^{\nu_{\mathrm{net}}+3}\le C_0d^{C_1}$ as soon as the piece count is at
most $md$, without any unspecified $d_*$. No theorem of this paper relies on that gate count.
\end{remark}

\begin{remark}[Lifted precision]\label{rem:transfer-lift}
Proposition~\ref{prop:prop14-param} applies directly at the lifted precision $q_*=q+8d$ used by the Walsh gadgets of
Section~\ref{sec:walsh} (with width $w_*=w+8d$ and halved lower band constants $(a_d/2,b_d,a_r/2,b_r)$ once $q_*\le2q$;
Lemma~\ref{lem:walsh-lift}).
\end{remark}
